\documentclass[answers,12pt,addpoints]{exam} 
\usepackage{import}

\import{C:/Users/prana/OneDrive/Desktop/MathNotes}{style.tex}

% Header
\newcommand{\name}{Pranav Tikkawar}
\newcommand{\course}{Spatial Statistics}
\newcommand{\assignment}{Homework 5.2}
\author{\name}
\title{\course \ - \assignment}

\begin{document}
\maketitle

\newpage
\begin{questions}
\question I am not super familiar with the concepts of bayesian inference, but I will try to understand it to the best of my ability. 

\question Why is it that when running some of these simulations like MCMC or Metropolis-Hastings, we get seemingly reasonable results, when the theory states otherwise. Is there some value in doing this, and can we use this to our advantage in some way. I mean that, the difference between something almost surely not being true in math is distinct from something being effectively impossible in practice.

\question In practice how do we know that our priors are good if there are scenarios where we can get good results even with bad priors? Is there some way to test the quality of our priors, or is it just a matter of trial and error?

\question You mention that there is not a good frequentist solution to the problem of the reliance of the different parameters to each other in the choice of model. Does one approach do it better than the other to ensure reliable results, or is it just a matter of being careful?

\question Is there a connection of the parameters of the power law covariance and matern model. Specifically in the smooth parameter. This may sound too ideal, but does knowing the parameters of one model give us any information about the parameters of the other model?

\question I do not understand how a change in slope with $m$ in the log-likelihood function implies no invertible tranformation of parameters. would it not require a non-continuous transformation of slope to imply that there is no invertible transformation of parameters? Also is there a systematic way the slope changes with $m$ that we can use to our advantage in some way?

\end{questions}
\end{document}